\chapter*{Jak korzystać z tej książki}
\phantomsection
\addcontentsline{toc}{chapter}{Jak korzystać z tej książki}
\markboth{JAK KORZYSTAĆ Z TEJ KSIĄŻKI}{Jak korzystać z tej książki}

\noindent
Tej książki nie da się przeczytać. Trzeba ją przerobić. To nie jest przenośnia
ani zachęta: format jest tak zbudowany, że bierne czytanie nie daje prawie
nic, i poczujesz to po czterech stronach.

\section*{Ramka}

Każdy program to strumień małych, ponumerowanych porcji zwanych
\emph{ramkami}. Prawie każda ramka kończy się pytaniem \dash{} o liczbę, słowo
albo wiersz rachunku. \textbf{Odpowiedź znajduje się na początku następnej
ramki.}

\begin{warning}
Zasłoń następną ramkę, zanim doczytasz bieżącą do końca. Wystarczy kartka
papieru. Zapisz swoją odpowiedź \dash{} na papierze, nie w głowie \dash{} i dopiero
wtedy odsłoń.

To jedyna zasada, która nie podlega negocjacji. Złam ją, a książka zdegeneruje
się do przeciętnego zbioru rozwiązanych przykładów i cała przewaga metody
zniknie. Będziesz miał ochotę ją złamać dokładnie w tych ramkach, w których
najbardziej się liczy \dash{} czyli w tych, których nie jesteś pewien.
\end{warning}

\noindent
Rząd kropek \blank{} oznacza, że brakuje słowa, liczby albo wyrażenia i masz
je uzupełnić.

\section*{Rusztowanie}

W obrębie jednego programu przykłady przechodzą przez cztery poziomy, w tej
kolejności, i ta zmiana jest zamierzona:

\begin{enumerate}[leftmargin=1.6em, itemsep=3pt]
  \item przykład rozwiązany w całości, z pokazanym rozumowaniem;
  \item przykład rozwiązany z lukami, które uzupełniasz;
  \item \emph{teraz jeden przykład do wykonania dla ciebie} \dash{} rozwiązujesz
        sam i sprawdzasz w następnej ramce;
  \item zadania na końcu programu, które robisz całkowicie samodzielnie.
\end{enumerate}

\noindent
Zanim dojdziesz do poziomu czwartego, wykonasz daną technikę trzy razy z coraz
mniejszą pomocą \dash{} dlatego zwykle wydaje się łatwa, a nie onieśmielająca. Ten
efekt jest właśnie celem konstrukcji.

\section*{Siedem części każdego programu}

\begin{center}
\small
\begin{tabularx}{\linewidth}{@{}lX@{}}
\toprule
\textbf{Efekty kształcenia} & Co będziesz umiał \dash{} wypisane, zanim
  zaczniesz, żebyś mógł ocenić program według jego własnej obietnicy. \\
\textbf{Quiz} & Tylko programy podstawowe. Rozwiąż go \emph{przed} programem:
  mówi, czy w ogóle potrzebujesz tego programu, a jeśli tylko fragmentu, to
  którego. Każde pytanie wskazuje ramki, które go uczą. \\
\textbf{Ramki} & Sam program. Podsumowanie i Zadania sprawdzające
  również są numerowanymi ramkami i noszą ostatnie numery w programie. \\
\textbf{Podsumowanie} & Każda pozycja opatrzona numerem ramki, z której
  pochodzi, \mbox{jak [17]}. Podsumowanie jest więc również drogą powrotną. \\
\textbf{Czy potrafisz?} & Te same efekty kształcenia, jeden do jednego, do
  samooceny w skali od 1 do 5. \\
\textbf{Zadania sprawdzające} & Dokładnie to, czego uczył program, i w tej
  samej kolejności. Bez pułapek. \\
\textbf{Dalsze wyzwania} & Utrwalenie \dash{} i więcej, niż będziesz chciał. \\
\bottomrule
\end{tabularx}
\end{center}

\section*{Pętla, która naprawdę działa}

\begin{center}
\begin{tcolorbox}[enhanced, sharp corners, width=0.94\linewidth,
  colback=mlight, colframe=mblue, boxrule=0pt, leftrule=3pt,
  left=10pt, right=10pt, top=8pt, bottom=8pt]
przeczytaj efekty $\rightarrow$ rozwiąż \textbf{Quiz} $\rightarrow$ przerób
program $\rightarrow$ \textbf{Czy potrafisz?} $\rightarrow$ \textbf{rozwiąż
Quiz jeszcze raz} $\rightarrow$ Zadania sprawdzające $\rightarrow$ Dalsze
wyzwania
\end{tcolorbox}
\end{center}

\noindent
Większość ludzi pomija powtórne rozwiązanie Quizu, a to właśnie ono niesie
najwięcej informacji. Quiz jest testem diagnostycznym na wejściu i testem
sprawdzającym na wyjściu; różnica między wynikami to jedyna uczciwa miara
tego, co program ci dał \dash{} i bywa wyraźnie mniejsza niż poczucie
zrozumienia.

\section*{Jeśli masz już wykształcenie techniczne}

Nie czytaj Części~I po kolei. Nie zakłada ona niczego, i to celowo; jeśli
umiesz zróżniczkować iloczyn, będzie ci się dłużyła.

\textbf{Rozwiąż trzynaście Quizów podstawowych i nic poza tym.} Zajmą ci
popołudnie. Wchodź w program podstawowy tylko tam, gdzie Quiz wypadnie źle, i
wchodź od ramek wskazanych przez nietrafione pytanie, a nie od ramki~1.
Następnie zacznij od Programu~1 i czytaj główne części po kolei \dash{} nie są
niezależne, a lista zależności we wstępie mówi, który program czego wymaga.

Dwa programy podstawowe warto przerobić nawet przy dobrym wyniku Quizu:
Program~F3 o logarytmach, bo prawie wszystko, co numeryczne w tej książce,
dzieje się w skali logarytmicznej, a powody nie są tymi ze szkoły; oraz
Program~F12 o regule łańcuchowej, bo propagacja wsteczna \emph{jest} regułą
łańcuchową i niczym więcej, a książka opiera się na tym zdaniu przez czterysta
stron.

\section*{Co ta książka twierdzi, a czego nie}

\begin{note}
Programy Strouda zostały przed publikacją zweryfikowane do poziomu
\textbf{80/80}: co najmniej 80 procent uczniów uzyskiwało co najmniej 80
procent. Wart uwagi wynik to jednak nie wzrost średniej, lecz zmniejszenie
\emph{rozstępu} \dash{} metoda podnosi dół rozkładu, a nie górę.

\textbf{Tej książki nie przetestowano na ani jednym czytelniku.} Pożycza metodę
i nie ma prawa pożyczyć dowodów. Każde jej budowanie wypisuje to jako zaległość
i będzie wypisywać, dopóki ktoś nie przeprowadzi badania.
\end{note}

\begin{warning}
\textbf{A ogólne dowody na skuteczność nauczania programowanego nie są dobre.}
Metaanaliza Kulika, Cohena i Ebelinga z 1982 roku, obejmująca ten format w
szkolnictwie średnim, wykazała, że zwykle nie podnosił on wyników ani nie
sprawiał, że uczniowie odnosili się do przedmiotu przychylniej. Małe kroki i
liniowa sekwencja nie okazały się niezbędne; natychmiastowe wzmocnienie nie
okazało się kluczowe.

To jest wynik przeciw temu formatowi jako całości i jest tu przytoczony, a nie
pominięty.

Powód, dla którego ta książka mimo to z formatu korzysta, jest węższy i lepiej
udokumentowany niż \enquote{nauczanie programowane działa}. Jest nim
\emph{przypominanie z pamięci}: przy stałym łącznym czasie nauki zmuszenie do
odtworzenia czegoś daje wyraźnie lepsze zapamiętanie długoterminowe niż
ponowne przeczytanie. I jest nim \emph{generowanie błędu}: wywołanie błędnej
odpowiedzi, a potem jej skorygowanie, bije naukę tego samego materiału bez
błędu \dash{} a korzyść jest największa dokładnie wtedy, gdy uczący się był
pewny swego. Te dwa wyniki są tym, po co są ramka i pułapka. Nie są obroną
małych kroków dla nich samych i ta książka takiej obrony nie oferuje.
\end{warning}

\begin{warning}
\textbf{W trakcie będzie się to wydawało gorsze niż czytanie.} To też nie jest
figura retoryczna. W badaniach nad przypominaniem z pamięci przewaga nie
istnieje w teście natychmiastowym, a pojawia się po dwóch dniach i po tygodniu;
badani konsekwentnie oceniają jako skuteczniejszą tę metodę, z której zapamiętują
\emph{mniej}, bo płynność czytania czuje się jak uczenie, a wysiłek jak porażka.

Jeśli porzucisz tę książkę, bo wydaje ci się nieefektywna, porzucisz ją z
powodu, dla którego ona działa.
\end{warning}
